
This paper documents a complete implementation failure in an automated research pipeline (YouRA) that reveals a missing validation checkpoint between experiment design and implementation phases. An investigation into LoRA-MoE coordination for multi-task learning produced zero experimental results because the selected model (Mixtral-8x7B, 47 billion parameters) required approximately 489GB VRAM, exceeding available capacity (5× H100 GPUs = 475GB total). Implementation proceeded to completion (29 files, 10 test suites, 100\% task completion, all quality checks passing) before discovering this constraint violation at execution time. The root cause is the absence of a computational feasibility gate in this pipeline: no validation checkpoint exists between Phase 2C (experiment design) and Phase 3-4 (implementation planning and coding) to estimate resource requirements against available hardware. We analyze this failure case, quantify implementation time invested before discovering infeasibility (10-16 hours), and propose a Phase 2C.5 feasibility gate design that estimates total memory requirements (model + optimizer + gradients + activations + overhead) with validation against hardware. Retrospective application shows the proposed gate would have correctly flagged this specific configuration as infeasible, requiring 489GB against 475GB capacity. This represents a meta-contribution: identifying a gap in one automated pipeline, quantifying its cost in this case, and proposing a solution that may benefit similar systems, subject to validation across multiple projects and diverse configurations.
